\documentclass[11pt]{article}
\usepackage{amsmath,amssymb,geometry}
\geometry{margin=22mm}
\begin{document}

\section*{P408: locker delivery, empty-flight counts and executable schedule}
Source: graph model (3)--(6) and the scheduling argument on pp.4--8.
The extracted file incorrectly turned illustrative duplicate tasks,
proof-specific sets/endpoints and numerical examples into universal
constraints. This independent repair removes those extraction artifacts
and states the executed graph problem and its operational certificate.

\paragraph{Data.}
$V$ is the supplied set of six locker sites.
Each task $t\in T$ is a distinct required loaded flight from $o_t$ to $d_t$,
possibly sharing its OD with another task, with release time $b_t$.
Initial parked drones are $s_v$, capacity is $c_v\ge s_v$, and
$r_v=c_v-s_v$. Distances $\lambda_{ij}$ are the supplied kilometres.
Candidate empty flights are $E=\{(i,j):i,j\in V,\ i\ne j\}$.
Define
\[
d_v^F=\#\{t:o_t=v\}-\#\{t:d_t=v\},\qquad
\mu_v^{ij}={\bf1}\{i=v\}-{\bf1}\{j=v\}.
\]
Use this out-minus-in convention throughout. It makes final parking
$s_v-d_v^F-\sum_{ij}\mu_v^{ij}x_{ij}$.

\paragraph{Integer graph optimization.}
Let $V_0=\{v:s_v=0\}$, and let $\mathcal C$ contain every nonempty
$W\subseteq V_0$ which touches at least one task and is not crossed by a
fixed task arc. The exact graph model executed is
\begin{align}
\min\quad & z_E=\sum_{(i,j)\in E}\lambda_{ij}x_{ij},\\
s_v-c_v \le\quad &d_v^F+\sum_{(i,j)\in E}\mu_v^{ij}x_{ij}\le s_v
 &&(v\in V),\\
&\sum_{\substack{(i,j)\in E\\{\bf1}\{i\in W\}\ne{\bf1}\{j\in W\}}}
x_{ij}\ge1 &&(W\in\mathcal C),\\
&x_{ij}\in\mathbb Z_{\ge0} &&((i,j)\in E).
\end{align}
All inactive subsets are enumerated for this six-site instance.
Every operationally feasible solution satisfies these degree and connection
conditions, so the graph optimum is also a lower bound on the operational
empty-flight distance.

\paragraph{Full operational feasibility and certificate.}
An operational solution must execute every labelled task once, together with
$x_{ij}$ empty flights of each OD. Each flight is assigned to a drone at its
origin; the drone may not overlap another flight, departures may not precede
release/loading readiness, and the number of parked drones at any site may
never exceed $c_v$ or fall below zero. The current algorithm searches a serial
execution of the optimal multiset. For a serial order $e_1,\ldots,e_m$, define
$q_v^0=s_v$ and impose for every step $h$, with $e_h=(i_h,j_h)$,
\[
q_{i_h}^{h-1}\ge1,\quad q_{j_h}^{h-1}<c_{j_h},\qquad
q_v^h=q_v^{h-1}-{\bf1}\{v=i_h\}+{\bf1}\{v=j_h\}.
\]
Assign an actually present drone at $i_h$. If loading/unloading takes
$\ell$ seconds for a loaded task and zero for an empty leg, and speed is
$v_{\rm fly}>0$, the returned timing is
\[
a_h=\max\{C_{h-1},b_{e_h}\},\quad
D_h=a_h+\ell_h,\quad
A_h=D_h+\frac{1000\lambda_{i_hj_h}}{v_{\rm fly}},\quad
C_h=A_h+\ell_h,\quad C_0=0.
\]
For empty legs $b_{e_h}=0$. The objective is determined by the stated routing costs.
Serial execution is a method for constructing feasible schedules. The existing instance has a feasible serial
certificate attaining the graph lower bound, which proves the same optimal
distance even for the larger operational scheduling domain.

The solver reports $z_E$ as \texttt{objective\_value}; total distance is
\[
z_{\rm total}=z_E+\sum_{t\in T}\lambda_{o_td_t},
\]
whose second term is constant. Both values are kept separately.
Code lines 6--20 implement the graph model; 23--42 construct a schedule,
and 43--56 independently replay parking, drone continuity, every task,
degree constraints and objective.

\end{document}
